\documentclass[11pt,a4paper]{article}
\usepackage[T1]{fontenc}
\usepackage[spanish,es-noshorthands]{babel}
\usepackage{lmodern,amsmath,amssymb,amsthm,mathtools}
\usepackage{graphicx,xcolor,enumitem,booktabs,array}
\usepackage[left=22mm,right=22mm,top=23mm,bottom=22mm,headheight=15pt]{geometry}
\usepackage{fancyhdr}
\usepackage[hidelinks]{hyperref}
\definecolor{uni}{RGB}{120,16,24}
\pagestyle{fancy}\fancyhf{}
\fancyhead[L]{\small\textcolor{uni}{UNI \textperiodcentered{} Escuela Profesional de Matemática}}
\fancyhead[R]{\small CM4H1 \textperiodcentered{} 2026-2}
\fancyfoot[L]{\footnotesize Oswaldo Velásquez Castañón}
\fancyfoot[R]{\footnotesize\thepage}
\renewcommand{\headrulewidth}{0.35pt}
\setlength{\parindent}{0pt}
\setlength{\parskip}{5pt}
\setlength{\emergencystretch}{2em}
\setlist[enumerate]{label=\alph*),leftmargin=7mm,itemsep=4pt,topsep=4pt}
\newcommand{\PP}{\mathbb P}
\newcommand{\EE}{\mathbb E}
\newcommand{\RR}{\mathbb R}
\newcommand{\Var}{\operatorname{Var}}
\newcommand{\Cov}{\operatorname{Cov}}
\newcommand{\Normal}{\mathcal N}
\newcommand{\puntos}[1]{\unskip\nobreak\hfill\mbox{\textnormal{\small[#1~\ifnum#1=1 punto\else puntos\fi]}}}
\newcommand{\problema}[3]{\par\vspace{4pt}{\large\bfseries\color{uni}#1. #2}\puntos{#3}\par\vspace{2pt}}
\newcommand{\encabezado}[2]{%
  \begin{minipage}{24mm}\includegraphics[width=20mm]{logo.png}\end{minipage}%
  \begin{minipage}{\dimexpr\textwidth-25mm\relax}
    {\LARGE\bfseries\color{uni}#1\par}
    \vspace{3pt}{\large Semana 7 \textperiodcentered{} Teorema central del límite\par}
    \vspace{3pt}{\small Teoría de la Probabilidad \textperiodcentered{} #2\par}
  \end{minipage}\par\vspace{5pt}\textcolor{uni}{\rule{\textwidth}{0.6pt}}\par}
\newcommand{\criterio}[1]{\par\begingroup\small\textbf{Criterio de calificación.} #1\par\endgroup}
\hypersetup{pdfauthor={Oswaldo Velásquez Castañón},pdftitle={Trabajo entregable \textperiodcentered{} Semana 7 \textperiodcentered{} TCL}}

\begin{document}
\encabezado{Trabajo entregable}{Actividad individual \textperiodcentered{} 20 puntos}
\textbf{Apellidos y nombres:} \hrulefill\quad \textbf{Código:} \rule{27mm}{0.4pt}

\textbf{Propósito.} Desarrollar una prueba alternativa del TCL y resolver problemas de transferencia que requieren argumentos distintos de los desarrollados en las notas y la práctica de la semana 7.

\textbf{Entrega.} Un PDF legible con cinco soluciones numeradas. Justifique cada paso al límite, indique las hipótesis utilizadas y distinga resultados exactos de aproximaciones. Dedicación orientativa: 6 a 8 horas. Puede usar calculadora o software para los valores numéricos, indicando cuál.

\textbf{Herramientas.} Puede citar resultados de clase sin reproducir sus pruebas. En el problema 1 se debe desarrollar el reemplazo gaussiano: no se admite sustituirlo por la expansión de la función característica y el límite de potencias de las notas. En los problemas 2--5 puede aplicar el TCL. No se utiliza Slutsky. Escriba $\Phi$ para la distribución normal estándar y $z_{0.975}\approx1.96$.

\problema{1}{Una prueba por reemplazo gaussiano}{6}
Sean $Y_k$ i.i.d. con $\EE Y_1=0$ y $\EE Y_1^2=1$. En un espacio producto tome $G_k$ i.i.d. $\Normal(0,1)$, independientes de todas las $Y_k$. Defina
\[
H_{n,k}=\frac{Y_1+\cdots+Y_k+G_{k+1}+\cdots+G_n}{\sqrt n},
\qquad 0\le k\le n,
\]
con sumas vacías iguales a cero. Sea $f\in C^2(\RR)$, con $f,f',f''$ acotadas y $f''$ uniformemente continua, y ponga
\[
\omega_f(\delta)=
\sup_{|x-y|\le\delta}|f''(x)-f''(y)|,\qquad \delta\ge0.
\]
\begin{enumerate}
\item Pruebe $\omega_f(\delta)\to0$ cuando $\delta\downarrow0$ y
$\omega_f(\delta)\le2\|f''\|_\infty$. Obtenga, por Taylor con resto integral,
\[
\left|f(x+h)-f(x)-hf'(x)-\frac{h^2}{2}f''(x)\right|
\le\frac{h^2}{2}\omega_f(|h|).
\] \puntos{1}
\item Para cada reemplazo de $G_k$ por $Y_k$, separe la suma que no contiene la variable reemplazada. Justifique que los términos lineales y cuadráticos tienen las mismas esperanzas. Escriba la identidad telescópica que compara $\EE f(H_{n,n})$ con $\EE f(H_{n,0})$. \puntos{2}
\item Deduzca una cota para esa diferencia en términos de
$\EE[Y_1^2\omega_f(|Y_1|/\sqrt n)]$ y
$\EE[G_1^2\omega_f(|G_1|/\sqrt n)]$.
Pruebe que tiende a cero usando solo segundos momentos. \puntos{2}
\item Aplique lo anterior a $f(x)=\cos(tx)$ y $f(x)=\sin(tx)$ para cada $t$ fijo. Use la ley de $H_{n,0}$ y Lévy para concluir el TCL de $n^{-1/2}\sum_{k=1}^nY_k$. Exprese el resultado para $X_k=m+\sigma Y_k$, $\sigma>0$. \puntos{1}
\end{enumerate}

\newpage
{\Large\bfseries\color{uni}Agrupación e inversión de sumas}\par
\problema{2}{Pesos alternados y una observación adicional}{4}
Sean $X_k$ i.i.d. con media cero y varianza uno. Defina
\[
V_n=\frac{\sum_{k=1}^n(X_{2k-1}+2X_{2k})}{\sqrt{5n}},
\qquad
W_n=V_n+\frac{X_{2n+1}}{\sqrt{5n}}.
\]
\begin{enumerate}
\item Calcule las medias y varianzas de $V_n$ y $W_n$. Explique por qué los sumandos individuales $X_{2k-1}$ y $2X_{2k}$ no tienen la misma ley. \puntos{1}
\item Proponga una agrupación que permita aplicar directamente el TCL i.i.d. y determine la ley límite de $V_n$. Justifique la independencia de los bloques. \puntos{1}
\item Pruebe que $W_n$ tiene el mismo límite usando funciones características y la independencia de la observación adicional. No se admite citar un teorema sobre perturbaciones que convergen a cero. \puntos{1}
\item Si el numerador de $V_n$ se divide entre $\sqrt{2n}$, ¿cuál es la ley límite? Justifique su respuesta e identifique la diferencia entre contar observaciones y sumar sus varianzas. \puntos{1}
\end{enumerate}

\problema{3}{Un conteo obtenido por inversión}{3}
Sean $E_k$ i.i.d. exponenciales de tasa $\lambda>0$, con densidad
$\lambda e^{-\lambda x}\mathbf1_{(0,\infty)}(x)$. Ponga
\[
S_0=0,\qquad S_j=\sum_{k=1}^jE_k,\qquad
N(t)=\max\{j\ge0:S_j\le t\},\qquad t\ge0.
\]
Demuestre
\[
\frac{N(t)-\lambda t}{\sqrt{\lambda t}}\Rightarrow\Normal(0,1)
\qquad(t\to\infty)
\]
sin utilizar la ley Poisson de $N(t)$ ni calcular su función característica.
\begin{enumerate}
\item Justifique que $N(t)$ es finito casi seguramente (puede usar la ley fuerte). Para $x\in\RR$ fijo, tome
$k_t=\lfloor\lambda t+x\sqrt{\lambda t}\rfloor$ y pruebe, para $t$ grande,
$\{N(t)\le k_t\}=\{S_{k_t+1}>t\}$. \puntos{1}
\item Calcule los momentos de $E_1$, aplique el TCL a $S_j$ y pruebe que
\[
a_t=\frac{\lambda t-(k_t+1)}{\sqrt{k_t+1}}\longrightarrow-x.
\] \puntos{1}
\item Justifique el paso al límite con el umbral móvil $a_t$ acotándolo entre $-x-\eta$ y $-x+\eta$, para $\eta>0$. Concluya el resultado del conteo. \puntos{1}
\end{enumerate}

\newpage
{\Large\bfseries\color{uni}Diseño de simulación y alcance del TCL}\par
\problema{4}{Montecarlo con parejas antitéticas}{4}
Para estimar $I=\int_0^1e^u\,du$, se comparan dos procedimientos que usan $2n$ evaluaciones de la función exponencial. Todas las variables $U_j$ y $V_k$ son uniformes en $[0,1]$ y mutuamente independientes:
\[
M_{2n}=\frac1{2n}\sum_{j=1}^{2n}e^{U_j},\qquad
A_n=\frac1n\sum_{k=1}^n Q_k,\quad
Q_k=\frac{e^{V_k}+e^{1-V_k}}{2}.
\]
\begin{enumerate}
\item Calcule $I$ y $v=\Var(e^{U_1})$. \puntos{1}
\item Calcule $\EE Q_1$, $c=\Cov(e^{V_1},e^{1-V_1})$ y $w=\Var(Q_1)$. Distinga la dependencia dentro de una pareja de la independencia entre parejas. \puntos{1}
\item Formule los TCL de $M_{2n}$ y $A_n$ con varianzas explícitas. Compare sus varianzas y los semianchos normales nominales al $95\%$ con el mismo costo $2n$. Evalúe la razón $\Var(A_n)/\Var(M_{2n})$. \puntos{1}
\item Obtenga el menor entero $n$ que satisface, para cada procedimiento, la regla normal de semiancho a lo sumo $0.01$. Use $z=1.96$ para esta regla. Informe también el número de evaluaciones y explique el alcance aproximado del cálculo. \puntos{1}
\end{enumerate}

\problema{5}{Eventos raros: un límite que no es normal}{3}
Para cada $n\ge2$, sean $I_{n,1},\ldots,I_{n,n}$ independientes, Bernoulli de parámetro $1/n$. No se supone relación entre filas. Defina
\[
B_n=\sum_{k=1}^nI_{n,k},\qquad
Z_n=\frac{B_n-\EE B_n}{\sqrt{\Var(B_n)}}.
\]
\begin{enumerate}
\item Calcule $\EE B_n$, $\Var(B_n)$ y la expresión explícita de $Z_n$. ¿Por qué no se puede aplicar directamente el TCL i.i.d. de las notas a estas filas? \puntos{1}
\item Calcule $\varphi_{Z_n}(t)$ y determine su límite para cada $t$. Identifique la ley límite mediante Lévy. Puede usar el lema de potencias complejas de las notas, sin demostrarlo nuevamente. \puntos{1}
\item Calcule $\lim_{n\to\infty}\PP(Z_n\le-1/2)$ y compárelo con $\Phi(-1/2)$. Explique qué refuta este ejemplo respecto de la afirmación \emph{un número grande de sumandos independientes siempre produce una aproximación normal}. \puntos{1}
\end{enumerate}

\vfill
{\small\textbf{Criterio de evaluación.} Se califica el desarrollo de los argumentos nuevos pedidos en cada inciso. Los resultados de las notas sirven como herramientas; su reproducción no reemplaza las soluciones.}
\end{document}
